\noindent\textbf{Resultado.} Además de $y\equiv0$, para cada $a\ge0$ y $\sigma\in\{-1,1\}$:
\[y(x)=\begin{cases}0,&x\le a,\\ \sigma[\tfrac23(x-a)]^{3/2},&x>a.\end{cases}\]

\noindent\textit{Desarrollo.} El campo $\sqrt[3]{y}$ es continuo, por lo que existe solución local. No es localmente Lipschitz en $y=0$. La función nula y las dos ramas que salen en $a=0$ ya prueban que no hay unicidad. Para obtener todas las ramas no nulas, derivamos $|y|^{2/3}$: en cualquier intervalo donde $y\ne0$, su derivada es $2/3$. Así $|y|^{2/3}=\frac23(x-a)$, necesariamente con $x>a$, y el signo de $y$ es constante. Una rama no nula no puede regresar a cero hacia adelante; antes de salir del cero sólo puede permanecer en él. Como $y(0)=0$, la salida debe satisfacer $a\ge0$. En el empalme, tanto $y$ como $y'$ tienden a cero, y la fórmula por tramos es $C^1$ y cumple la EDO en todo punto. Permitir cualquier tiempo de espera, o no salir nunca, completa la familia.
